% !TeX root = part02.tex

\documentclass[../final.tex]{subfiles}

\begin{document}

% ============================================================
% Лекция 08.09
% Решение нелинейных уравнений
% ============================================================

\chapter{Решение нелинейных уравнений}

\rsubtitle{Методы дихотомии, простых итераций и Ньютона}

\captionsetup{
    font=footnotesize,
    labelfont=bf,
    skip=5pt
}

% Общий стиль рисунков лекции.
\tikzset{
    ovfaxis/.style={
        -{Stealth[length=2.2mm]},
        white,
        line width=0.85pt
    },
    ovfcurve/.style={
        rc4,
        line width=1.55pt
    },
    ovfcurvealt/.style={
        white!80,
        line width=1.05pt
    },
    ovfguide/.style={
        white!45,
        dashed,
        line width=0.70pt
    },
    ovfiter/.style={
        rc2,
        -{Stealth[length=1.8mm]},
        line width=1.0pt
    },
    ovftangent/.style={
        rc2,
        line width=1.0pt
    },
    ovfpoint/.style={
        circle,
        fill=rc3,
        draw=none,
        inner sep=1.7pt
    }
}

% Функция, используемая на рисунках метода Ньютона и секущих.
\pgfmathdeclarefunction{ovff}{1}{%
    \pgfmathparse{0.25*(#1-1.2)^3 + 0.8*(#1-1.2)}%
}
\pgfmathdeclarefunction{ovffp}{1}{%
    \pgfmathparse{0.75*(#1-1.2)^2 + 0.8}%
}


% ============================================================
\section{Постановка задачи}
% ============================================================

Будем рассматривать численное решение нелинейного уравнения
одной вещественной переменной
%
\begin{equation*}
    f(x)=0.
\end{equation*}

В общем случае такое уравнение не удаётся решить аналитически.
Простейший пример --- уравнение
%
\begin{equation*}
    x=\cos x,
\end{equation*}
%
которое можно переписать в стандартном виде
%
\begin{equation*}
    x-\cos x=0.
\end{equation*}

Геометрически его корень определяется как точка пересечения графиков
$y=x$ и $y=\cos x$.

\begin{center}
\begin{tikzpicture}[x=2.45cm,y=2.45cm]
    \draw[ovfaxis] (-0.35,0) -- (2.15,0)
        node[below left=1pt,text=white] {$x$};
    \draw[ovfaxis] (0,-0.55) -- (0,1.45)
        node[left=1pt,text=white] {$y$};

    \draw[ovfcurvealt,domain=-0.25:1.35,samples=2]
        plot (\x,{\x});
    \draw[ovfcurve,domain=-0.25:2.0,samples=100,smooth]
        plot (\x,{cos(deg(\x))});

    \fill[rc3] (0.7391,0.7391) circle (1.6pt);

    \node[text=white!85,anchor=south west]
        at (1.15,1.15) {$y=x$};
    \node[text=rc4,anchor=south west]
        at (1.25,0.28) {$y=\cos x$};
    \node[text=white!80,anchor=north west]
        at (0.78,0.70) {$x^*$};
\end{tikzpicture}

\captionof{figure}{Геометрическая интерпретация уравнения $x=\cos x$.}
\end{center}

Аналогично уравнение
%
\begin{equation*}
    x=\tan x
\end{equation*}
%
можно записать как
%
\begin{equation*}
    x-\tan x=0.
\end{equation*}

В отличие от предыдущего примера, такое уравнение имеет несколько
вещественных корней.

\begin{center}
\begin{tikzpicture}[x=0.75cm,y=0.75cm]
    \clip (-5.3,-4.2) rectangle (5.3,4.2);

    \draw[ovfaxis] (-5.15,0) -- (5.15,0)
        node[below,text=white] {$x$};
    \draw[ovfaxis] (0,-4.0) -- (0,4.0)
        node[left,text=white] {$y$};

    \draw[ovfcurvealt,domain=-4.1:4.1,samples=2]
        plot (\x,{\x});

    \foreach \a/\b in {
        -4.68/-1.62,
        -1.52/1.52,
        1.62/4.68
    }{
        \draw[ovfcurve,domain=\a:\b,samples=120,smooth]
            plot (\x,{tan(deg(\x))});
    }

    \foreach \x in {-4.4934,0,4.4934}{
        \fill[rc3] (\x,\x) circle (1.5pt);
    }

    \node[text=white!85,anchor=south west]
        at (2.7,2.7) {$y=x$};
    \node[text=rc4,anchor=south east]
        at (1.15,3.2) {$y=\tan x$};
\end{tikzpicture}

\captionof{figure}{Несколько ветвей $y=\tan x$ и прямая $y=x$.}
\end{center}

Таким образом, прежде чем применять численный метод, необходимо
выбрать интересующий нас корень и локализовать его на некотором
отрезке.


% ============================================================
\section{Изоляция корня}
% ============================================================

Пусть требуется найти корень $x^*$ уравнения
%
\begin{equation*}
    f(x)=0.
\end{equation*}

Найдём интервал
%
\begin{equation*}
    [a,b],
\end{equation*}
%
содержащий интересующий нас корень. Будем считать, что на этом
интервале выполняются следующие условия:

\begin{enumerate}[label=\arabic*)]
    \item на $[a,b]$ находится только один корень $x^*$;
    \item функция $f(x)$ монотонна на $[a,b]$;
    \item функция $f(x)$ непрерывна на $[a,b]$;
    \item значения функции на концах имеют разные знаки либо один
    из концов сам является корнем:
    %
    \begin{equation*}
        f(a)f(b)\leq0.
    \end{equation*}
\end{enumerate}

Для непрерывной функции последнее условие гарантирует существование
хотя бы одного корня внутри отрезка. Дополнительное условие
единственности позволяет говорить именно об изолированном корне.

\begin{center}
\begin{tikzpicture}[x=1.5cm,y=1.0cm]
    \draw[ovfaxis] (-0.15,0) -- (5.8,0)
        node[below,text=white] {$x$};
    \draw[ovfaxis] (0,-2.3) -- (0,2.7)
        node[left,text=white] {$y$};

    \draw[ovfcurve,domain=0.15:5.45,samples=120,smooth]
        plot (\x,{0.34*(\x-1.5)*(\x-4.5)});

    \draw[ovfguide] (0.70,0) -- (0.70,{0.34*(0.70-1.5)*(0.70-4.5)});
    \draw[ovfguide] (2.40,0) -- (2.40,{0.34*(2.40-1.5)*(2.40-4.5)});

    \fill[rc3] (1.5,0) circle (1.7pt);
    \fill[rc3] (4.5,0) circle (1.7pt);

    \node[text=white,below] at (0.70,0) {$a$};
    \node[text=white,below] at (2.40,0) {$b$};
    \node[text=white!85,above] at (1.5,0) {$x_1^*$};
    \node[text=white!85,above] at (4.5,0) {$x_2^*$};
    \node[text=rc4,anchor=south west] at (4.25,1.0) {$y=f(x)$};
\end{tikzpicture}

\captionof{figure}{Изоляция одного корня на отрезке $[a,b]$.}
\end{center}


% ============================================================
\section{Метод дихотомии}
% ============================================================

\subsection{Построение последовательности отрезков}

Пусть первоначально
%
\begin{equation*}
    a_0=a,
    \qquad
    b_0=b,
\end{equation*}
%
и
%
\begin{equation*}
    f(a_0)f(b_0)\leq0.
\end{equation*}

На первом шаге делим исходный отрезок пополам и вычисляем его
середину:
%
\begin{equation*}
    x_1=\frac{a_0+b_0}{2}.
\end{equation*}

Если
%
\begin{equation*}
    f(a_0)f(x_1)\leq0,
\end{equation*}
%
то выбираем левую половину:
%
\begin{equation*}
    a_1=a_0,
    \qquad
    b_1=x_1.
\end{equation*}

Если же
%
\begin{equation*}
    f(a_0)f(x_1)>0,
\end{equation*}
%
то корень находится в правой половине и полагаем
%
\begin{equation*}
    a_1=x_1,
    \qquad
    b_1=b_0.
\end{equation*}

В обоих случаях сохраняется условие
%
\begin{equation*}
    f(a_1)f(b_1)\leq0,
\end{equation*}
%
а ширина нового интервала в два раза меньше исходной:
%
\begin{equation*}
    b_1-a_1
    =
    \frac{b_0-a_0}{2}.
\end{equation*}

Введём обозначение
%
\begin{equation*}
    \delta_0=b_0-a_0.
\end{equation*}
%
Тогда
%
\begin{equation*}
    \delta_1=\frac{\delta_0}{2}.
\end{equation*}

\begin{center}
\begin{tikzpicture}[x=2.0cm,y=1.15cm]
    \draw[ovfaxis] (-0.15,0) -- (3.25,0)
        node[below,text=white] {$x$};
    \draw[ovfaxis] (0,-1.4) -- (0,2.3)
        node[left,text=white] {$y$};

    \draw[ovfcurve,domain=0.55:2.75,samples=100,smooth]
        plot (\x,{1.15*(\x*\x-2)});

    \pgfmathsetmacro{\aa}{1.0}
    \pgfmathsetmacro{\bb}{2.0}
    \pgfmathsetmacro{\xxone}{1.5}

    \draw[ovfguide] (\aa,0) -- (\aa,{1.15*(\aa*\aa-2)});
    \draw[ovfguide] (\bb,0) -- (\bb,{1.15*(\bb*\bb-2)});
    \draw[ovfguide] (\xxone,0) -- (\xxone,{1.15*(\xxone*\xxone-2)});

    \fill[rc3] ({sqrt(2)},0) circle (1.7pt);

    \node[text=white,below] at (\aa,0) {$a_0$};
    \node[text=white,below] at (\bb,0) {$b_0$};
    \node[text=white!85,below] at (\xxone,0) {$x_1$};
    \node[text=white!85,above left] at ({sqrt(2)},0) {$x^*$};
    \node[text=rc4,anchor=south west] at (2.15,1.35) {$y=f(x)$};
\end{tikzpicture}

\captionof{figure}{Первый шаг метода дихотомии.}
\end{center}

На втором шаге повторяем ту же процедуру уже для отрезка
$[a_1,b_1]$:
%
\begin{equation*}
    x_2=\frac{a_1+b_1}{2}.
\end{equation*}

После выбора нужной половины получаем
%
\begin{equation*}
    b_2-a_2
    =
    \delta_2
    =
    \frac{b_0-a_0}{2^2}
    =
    \frac{\delta_0}{2^2}.
\end{equation*}

Продолжая процесс, на $n$-м шаге имеем
%
\begin{equation*}
    x_n
    =
    \frac{a_{n-1}+b_{n-1}}{2},
\end{equation*}
%
а ширина нового интервала равна
%
\begin{align*}
    b_n-a_n
    &=
    \frac{b_{n-1}-a_{n-1}}{2}
    \\
    &=
    \frac{b_0-a_0}{2^n}
    \\
    &=
    \frac{\delta_0}{2^n}.
\end{align*}

Таким образом,
%
\begin{equation*}
    \delta_n
    =
    \frac{\delta_0}{2^n}.
\end{equation*}


\subsection{Число шагов}

Условием остановки итерационного процесса можно взять достижение
заданной точности
%
\begin{equation*}
    \delta_n\leq\varepsilon.
\end{equation*}

Тогда
%
\begin{equation*}
    \frac{\delta_0}{2^N}\leq\varepsilon,
\end{equation*}
%
откуда
%
\begin{equation*}
    2^N\geq\frac{\delta_0}{\varepsilon}.
\end{equation*}

Следовательно, необходимое число шагов оценивается как
%
\begin{equation*}
    N
    \approx
    \log_2\frac{\delta_0}{\varepsilon}
    =
    \frac{\ln(\delta_0/\varepsilon)}{\ln2}.
\end{equation*}

Если в качестве приближения к корню брать середину текущего
интервала, то максимальная абсолютная погрешность вдвое меньше его
ширины:
%
\begin{equation*}
    \left|x_n-x^*\right|
    \leq
    \frac{b_{n-1}-a_{n-1}}{2}.
\end{equation*}


\subsection{Сходимость метода}

Последовательность левых границ
%
\begin{equation*}
    \{a_n\}
\end{equation*}
%
монотонно возрастает и ограничена сверху. Последовательность правых
границ
%
\begin{equation*}
    \{b_n\}
\end{equation*}
%
монотонно убывает и ограничена снизу. Поэтому обе
последовательности имеют пределы:
%
\begin{equation*}
    a_n\to A,
    \qquad
    b_n\to B.
\end{equation*}

Но
%
\begin{equation*}
    b_n-a_n
    =
    \frac{\delta_0}{2^n}
    \underset{n\to\infty}{\longrightarrow}
    0,
\end{equation*}
%
следовательно,
%
\begin{equation*}
    A=B=C.
\end{equation*}

На каждом шаге сохраняется условие
%
\begin{equation*}
    f(a_n)f(b_n)\leq0.
\end{equation*}

Переходя к пределу и используя непрерывность функции $f(x)$,
получаем
%
\begin{equation*}
    f(C)f(C)\leq0.
\end{equation*}

То есть
%
\begin{equation*}
    [f(C)]^2\leq0.
\end{equation*}

Поскольку квадрат вещественного числа неотрицателен,
%
\begin{equation*}
    f(C)=0.
\end{equation*}

Следовательно,
%
\begin{equation*}
    C=x^*,
\end{equation*}
%
и
%
\begin{equation*}
    x_n
    =
    \frac{a_{n-1}+b_{n-1}}{2}
    \underset{n\to\infty}{\longrightarrow}
    x^*.
\end{equation*}

Таким образом, метод дихотомии обладает гарантированной сходимостью
для непрерывной функции при наличии изменения знака на концах
отрезка. Недостаток метода --- сравнительно невысокая скорость
сходимости.


% ============================================================
\section{Метод простых итераций}
% ============================================================

\subsection{Переход к уравнению на неподвижную точку}

Исходное уравнение
%
\begin{equation*}
    f(x)=0
\end{equation*}
%
перепишем в эквивалентном виде
%
\begin{equation*}
    x=\varphi(x).
\end{equation*}

Например, простейший вариант такого преобразования имеет вид
%
\begin{equation*}
    f(x)=0
    \quad\Longleftrightarrow\quad
    x+f(x)=x,
\end{equation*}
%
то есть
%
\begin{equation*}
    \varphi(x)=x+f(x).
\end{equation*}

После этого строим итерационный процесс
%
\begin{equation*}
    x_{n+1}=\varphi(x_n),
    \qquad
    n=0,1,2,\ldots
\end{equation*}

Выбираем начальное приближение $x_0$ и последовательно вычисляем
%
\begin{equation*}
    x_1=\varphi(x_0),
    \qquad
    x_2=\varphi(x_1),
    \qquad
    x_3=\varphi(x_2),
    \ldots
\end{equation*}

\begin{center}
\begin{tikzpicture}[x=3.0cm,y=3.0cm]
    \pgfmathsetmacro{\xzero}{1.70}
    \pgfmathsetmacro{\xone}{0.6+0.3*cos(deg(\xzero))}
    \pgfmathsetmacro{\xtwo}{0.6+0.3*cos(deg(\xone))}
    \pgfmathsetmacro{\xthree}{0.6+0.3*cos(deg(\xtwo))}

    \draw[ovfaxis] (-0.05,0) -- (1.95,0)
        node[below,text=white] {$x$};
    \draw[ovfaxis] (0,-0.05) -- (0,1.45)
        node[left,text=white] {$y$};

    \draw[ovfcurvealt,domain=0:1.45,samples=2]
        plot (\x,{\x});
    \draw[ovfcurve,domain=0:1.85,samples=100,smooth]
        plot (\x,{0.6+0.3*cos(deg(\x))});

    \draw[ovfiter] (\xzero,0) -- (\xzero,\xone);
    \draw[ovfiter] (\xzero,\xone) -- (\xone,\xone);
    \draw[ovfiter] (\xone,\xone) -- (\xone,\xtwo);
    \draw[ovfiter] (\xone,\xtwo) -- (\xtwo,\xtwo);
    \draw[ovfiter] (\xtwo,\xtwo) -- (\xtwo,\xthree);
    \draw[ovfiter] (\xtwo,\xthree) -- (\xthree,\xthree);

    \node[text=white,below] at (\xzero,0) {$x_0$};
    \node[text=white,below] at (\xone,0) {$x_1$};
    \node[text=white,below] at (\xtwo,0) {$x_2$};

    \node[text=white!85,anchor=south west]
        at (1.05,1.05) {$y=x$};
    \node[text=rc4,anchor=south west]
        at (1.05,0.63) {$y=\varphi(x)$};
\end{tikzpicture}

\captionof{figure}{Геометрическая схема метода простых итераций.}
\end{center}


\subsection{Условие сходимости}

Пусть $[a,b]$ --- интервал, содержащий корень $x^*$, причём
%
\begin{equation*}
    \varphi([a,b])\subset[a,b].
\end{equation*}

Пусть функция $\varphi(x)$ непрерывно дифференцируема на $[a,b]$ и
существует число $q<1$ такое, что
%
\begin{equation*}
    |\varphi'(x)|\leq q<1,
    \qquad
    x\in[a,b].
\end{equation*}

Тогда итерационный процесс
%
\begin{equation*}
    x_{n+1}=\varphi(x_n)
\end{equation*}
%
сходится к единственной неподвижной точке $x^*$ на этом интервале.

Действительно, для корня
%
\begin{equation*}
    x^*=\varphi(x^*).
\end{equation*}

Поэтому
%
\begin{align*}
    |x_n-x^*|
    &=
    |\varphi(x_{n-1})-\varphi(x^*)|.
\end{align*}

По теореме Лагранжа существует точка $\xi_{n-1}$ между
$x_{n-1}$ и $x^*$ такая, что
%
\begin{align*}
    |x_n-x^*|
    &=
    |\varphi'(\xi_{n-1})|
    |x_{n-1}-x^*|
    \\
    &\leq
    q|x_{n-1}-x^*|.
\end{align*}

Для предыдущего шага аналогично
%
\begin{equation*}
    |x_{n-1}-x^*|
    \leq
    q|x_{n-2}-x^*|.
\end{equation*}

Следовательно,
%
\begin{align*}
    |x_n-x^*|
    &\leq
    q^2|x_{n-2}-x^*|
    \\
    &\leq
    \ldots
    \\
    &\leq
    q^n|x_0-x^*|.
\end{align*}

Так как
%
\begin{equation*}
    0\leq q<1,
\end{equation*}
%
то
%
\begin{equation*}
    q^n\underset{n\to\infty}{\longrightarrow}0,
\end{equation*}
%
и поэтому
%
\begin{equation*}
    x_n\underset{n\to\infty}{\longrightarrow}x^*.
\end{equation*}


\subsection{Сходящийся и расходящийся процессы}

Условие
%
\begin{equation*}
    |\varphi'(x)|<1
\end{equation*}
%
соответствует сжимающему отображению: расстояние до неподвижной
точки уменьшается после каждой итерации. Если же в окрестности корня
%
\begin{equation*}
    |\varphi'(x^*)|>1,
\end{equation*}
%
то неподвижная точка становится отталкивающей и итерации обычно
расходятся.

\begin{center}
\begin{minipage}{0.46\linewidth}
\centering
\begin{tikzpicture}[x=2.1cm,y=2.1cm]
    \draw[ovfaxis] (-0.05,0) -- (2.0,0);
    \draw[ovfaxis] (0,-0.05) -- (0,2.0);

    \draw[ovfcurvealt] (0,0) -- (1.85,1.85);
    \draw[ovfcurve,domain=0:1.85,samples=2]
        plot (\x,{1+0.55*(\x-1)});

    \node[text=white!85,anchor=south west] at (1.35,1.35) {$y=x$};
    \node[text=rc4,anchor=north west] at (1.28,1.15) {$y=\varphi(x)$};
    \node[text=white,below] at (1,0) {$x^*$};
    \fill[rc3] (1,1) circle (1.5pt);
\end{tikzpicture}

{\footnotesize $|\varphi'(x^*)|<1$: корень притягивает итерации.}
\end{minipage}
\hfill
\begin{minipage}{0.46\linewidth}
\centering
\begin{tikzpicture}[x=2.1cm,y=2.1cm]
    \draw[ovfaxis] (-0.05,0) -- (2.0,0);
    \draw[ovfaxis] (0,-0.05) -- (0,2.0);

    \draw[ovfcurvealt] (0,0) -- (1.85,1.85);
    \draw[ovfcurve,domain=0.22:1.72,samples=2]
        plot (\x,{1+1.25*(\x-1)});

    \node[text=white!85,anchor=south west] at (1.35,1.35) {$y=x$};
    \node[text=rc4,anchor=south east] at (0.75,0.55) {$y=\varphi(x)$};
    \node[text=white,below] at (1,0) {$x^*$};
    \fill[rc3] (1,1) circle (1.5pt);
\end{tikzpicture}

{\footnotesize $|\varphi'(x^*)|>1$: корень отталкивает итерации.}
\end{minipage}
\end{center}

Ограничение $|\varphi'|<1$ является достаточным, но не необходимым
условием сходимости. В более общем случае достаточно, чтобы само
отображение уменьшало расстояние до искомой неподвижной точки.


\subsection{Скорость сходимости}

Пусть начальное приближение $x_0$ уже достаточно близко к корню.
Разложим $\varphi(x_0)$ около $x^*$:
%
\begin{equation*}
    \varphi(x_0)
    \approx
    \varphi(x^*)
    +
    \varphi'(x^*)(x_0-x^*).
\end{equation*}

Так как
%
\begin{equation*}
    \varphi(x^*)=x^*,
\end{equation*}
%
то
%
\begin{align*}
    x_1-x^*
    &\approx
    \varphi'(x^*)(x_0-x^*).
\end{align*}

Введём абсолютную погрешность
%
\begin{equation*}
    \delta_n=|x_n-x^*|.
\end{equation*}

Тогда
%
\begin{equation*}
    \delta_1
    \approx
    |\varphi'(x^*)|\,\delta_0.
\end{equation*}

На следующем шаге
%
\begin{equation*}
    \delta_2
    \approx
    |\varphi'(x^*)|\,\delta_1
    \approx
    |\varphi'(x^*)|^2\delta_0.
\end{equation*}

Продолжая,
%
\begin{equation*}
    \delta_n
    \approx
    |\varphi'(x^*)|^n\delta_0.
\end{equation*}

Если требуется достичь точности $\varepsilon$, то
%
\begin{equation*}
    \varepsilon
    \approx
    |\varphi'(x^*)|^N\delta_0.
\end{equation*}

Следовательно,
%
\begin{equation*}
    N
    \approx
    \frac{\ln(\delta_0/\varepsilon)}
    {-\ln|\varphi'(x^*)|}.
\end{equation*}

Сравнивая эту оценку с методом дихотомии, получаем, что метод
простых итераций сходится быстрее при
%
\begin{equation*}
    |\varphi'(x^*)|<\frac12.
\end{equation*}

Особенно быстрой будет сходимость при
%
\begin{equation*}
    \varphi'(x^*)=0,
\end{equation*}
%
однако в этом случае линейная оценка уже недостаточна и необходимо
учитывать следующие члены разложения Тейлора.


\subsection{Переход к обратной функции}

Если
%
\begin{equation*}
    |\varphi'(x)|>1,
\end{equation*}
%
то прямой итерационный процесс
%
\begin{equation*}
    x_{n+1}=\varphi(x_n)
\end{equation*}
%
может расходиться. Если обратную функцию удобно вычислять, можно
перейти от
%
\begin{equation*}
    x=\varphi(x)
\end{equation*}
%
к эквивалентному уравнению
%
\begin{equation*}
    x=\varphi^{-1}(x).
\end{equation*}

При этом
%
\begin{equation*}
    (\varphi^{-1})'(x)
    =
    \frac{1}{\varphi'(\varphi^{-1}(x))}.
\end{equation*}

В окрестности неподвижной точки
%
\begin{equation*}
    |(\varphi^{-1})'(x^*)|
    =
    \frac{1}{|\varphi'(x^*)|}
    <1,
\end{equation*}
%
так что итерации для обратной функции становятся сходящимися.


\subsection{Корректирующий множитель}

Вместо простейшего представления
%
\begin{equation*}
    x=x+f(x)
\end{equation*}
%
можно использовать эквивалентное уравнение
%
\begin{equation*}
    x=x-\lambda f(x),
    \qquad
    \lambda\neq0.
\end{equation*}

Определим
%
\begin{equation*}
    \widetilde\varphi(x)
    =
    x-\lambda f(x).
\end{equation*}

Тогда
%
\begin{equation*}
    \widetilde\varphi'(x)
    =
    1-\lambda f'(x).
\end{equation*}

Для сходимости достаточно потребовать
%
\begin{equation*}
    |1-\lambda f'(x)|<1.
\end{equation*}

Раскрывая это неравенство,
%
\begin{align*}
    -1
    &<
    1-\lambda f'(x)
    <
    1,
    \\
    -2
    &<
    -\lambda f'(x)
    <
    0,
\end{align*}
%
то есть
%
\begin{equation*}
    0<\lambda f'(x)<2.
\end{equation*}

Поэтому знак $\lambda$ следует выбирать совпадающим со знаком
$f'(x)$. Если производная не меняет знак на всём рабочем интервале,
то достаточное условие можно записать как
%
\begin{equation*}
    0<|\lambda|
    <
    \frac{2}
    {\displaystyle\max_{x\in[a,b]}|f'(x)|}.
\end{equation*}

При этом для ускорения сходимости значение $\lambda$ желательно
выбирать как можно ближе к
%
\begin{equation*}
    \frac{1}{f'(x^*)}.
\end{equation*}

Итерационный процесс принимает вид
%
\begin{equation*}
    x_{n+1}
    =
    x_n-\lambda f(x_n).
\end{equation*}

\begin{center}
\begin{tikzpicture}[x=1.8cm,y=1.0cm]
    \pgfmathsetmacro{\xn}{3.0}
    \pgfmathsetmacro{\fyn}{ovff(\xn)}
    \pgfmathsetmacro{\slope}{1.90}
    \pgfmathsetmacro{\xnext}{\xn-\fyn/\slope}

    \draw[ovfaxis] (-0.10,0) -- (3.65,0)
        node[below,text=white] {$x$};
    \draw[ovfaxis] (0,-1.3) -- (0,3.55)
        node[left,text=white] {$y$};

    \draw[ovfcurve,domain=0.15:3.25,samples=120,smooth]
        plot (\x,{ovff(\x)});

    \draw[ovftangent]
        (\xnext,0)
        --
        (\xn,\fyn);

    \draw[ovfguide] (\xn,0) -- (\xn,\fyn);
    \fill[rc3] (\xn,\fyn) circle (1.6pt);
    \fill[rc3] (1.2,0) circle (1.6pt);

    \node[text=white,below] at (\xn,0) {$x_n$};
    \node[text=white,below] at (\xnext,0) {$x_{n+1}$};
    \node[text=white!85,below] at (1.2,0) {$x^*$};
    \node[text=rc4,anchor=south west] at (2.45,1.65) {$y=f(x)$};

    \draw[white!65,line width=0.7pt]
        (\xnext+0.28,0)
        arc[start angle=0,end angle=58,radius=0.28];
    \node[text=white!75] at (\xnext+0.34,0.38) {$\alpha$};
\end{tikzpicture}

\captionof{figure}{Итерация с корректирующим множителем:
наклон прямой задаётся величиной $1/\lambda$.}
\end{center}

Если угол прямой с осью $x$ обозначить через $\alpha$, то
%
\begin{equation*}
    \tan\alpha=\frac{1}{\lambda}.
\end{equation*}

Именно эта геометрическая интерпретация приводит к следующему
ускорению метода: вместо прямой с постоянным наклоном будем на каждом
шаге использовать касательную к графику $y=f(x)$.


% ============================================================
\section{Метод Ньютона}
% ============================================================

\subsection{Формула метода}

Для касательной к графику $y=f(x)$ в точке $x_n$ имеем
%
\begin{equation*}
    \tan\alpha=f'(x_n).
\end{equation*}

Сравнивая это с
%
\begin{equation*}
    \tan\alpha=\frac{1}{\lambda},
\end{equation*}
%
выбираем
%
\begin{equation*}
    \lambda(x_n)=\frac{1}{f'(x_n)}.
\end{equation*}

Тогда итерационный процесс принимает вид
%
\begin{equation*}
    x_{n+1}
    =
    x_n
    -
    \frac{f(x_n)}{f'(x_n)}.
\end{equation*}

Это \rtxt{3}{метод Ньютона}, или метод Ньютона--Рафсона.

Геометрически на каждом шаге проводим касательную к графику
$y=f(x)$ в точке $(x_n,f(x_n))$. Точка пересечения этой касательной
с осью $x$ даёт следующее приближение $x_{n+1}$.

\begin{center}
\begin{tikzpicture}[x=1.75cm,y=0.95cm]
    \pgfmathsetmacro{\xzero}{3.0}
    \pgfmathsetmacro{\fzero}{ovff(\xzero)}
    \pgfmathsetmacro{\fpzero}{ovffp(\xzero)}
    \pgfmathsetmacro{\xone}{\xzero-\fzero/\fpzero}
    \pgfmathsetmacro{\fone}{ovff(\xone)}
    \pgfmathsetmacro{\fpone}{ovffp(\xone)}
    \pgfmathsetmacro{\xtwo}{\xone-\fone/\fpone}

    \draw[ovfaxis] (-0.10,0) -- (3.65,0)
        node[below,text=white] {$x$};
    \draw[ovfaxis] (0,-1.3) -- (0,3.65)
        node[left,text=white] {$y$};

    \draw[ovfcurve,domain=0.15:3.25,samples=120,smooth]
        plot (\x,{ovff(\x)});

    \draw[ovftangent,domain=\xone:3.20,samples=2]
        plot (\x,{\fzero+\fpzero*(\x-\xzero)});
    \draw[ovftangent,domain=\xtwo:2.55,samples=2]
        plot (\x,{\fone+\fpone*(\x-\xone)});

    \draw[ovfguide] (\xzero,0) -- (\xzero,\fzero);
    \draw[ovfguide] (\xone,0) -- (\xone,\fone);

    \fill[rc3] (\xzero,\fzero) circle (1.6pt);
    \fill[rc3] (\xone,\fone) circle (1.6pt);
    \fill[rc3] (1.2,0) circle (1.6pt);

    \node[text=white,below] at (\xzero,0) {$x_0$};
    \node[text=white,below] at (\xone,0) {$x_1$};
    \node[text=white,below] at (\xtwo,0) {$x_2$};
    \node[text=white!85,below] at (1.2,0) {$x^*$};
    \node[text=rc4,anchor=south west] at (2.45,1.60) {$y=f(x)$};
\end{tikzpicture}

\captionof{figure}{Геометрическая схема метода Ньютона.}
\end{center}


\subsection{Квадратичная сходимость}

Оценим погрешность вблизи простого корня $x^*$, для которого
%
\begin{equation*}
    f(x^*)=0,
    \qquad
    f'(x^*)\neq0.
\end{equation*}

Введём знаковую погрешность
%
\begin{equation*}
    R_n=x^*-x_n.
\end{equation*}

Тогда
%
\begin{equation*}
    x_n-x^*=-R_n.
\end{equation*}

Из формулы Ньютона
%
\begin{equation*}
    x_{n+1}
    =
    x_n
    -
    \frac{f(x_n)}{f'(x_n)}
\end{equation*}
%
получаем
%
\begin{align*}
    R_{n+1}
    &=
    x^*-x_{n+1}
    \\
    &=
    x^*-x_n
    +
    \frac{f(x_n)}{f'(x_n)}
    \\
    &=
    R_n
    +
    \frac{f(x_n)}{f'(x_n)}.
\end{align*}

Разложим функцию около корня:
%
\begin{align*}
    f(x_n)
    &\approx
    f(x^*)
    +
    f'(x^*)(x_n-x^*)
    +
    \frac12 f''(x^*)(x_n-x^*)^2
    \\
    &=
    -f'(x^*)R_n
    +
    \frac12 f''(x^*)R_n^2.
\end{align*}

Аналогично для производной
%
\begin{align*}
    f'(x_n)
    &\approx
    f'(x^*)
    +
    f''(x^*)(x_n-x^*)
    \\
    &=
    f'(x^*)
    -
    f''(x^*)R_n.
\end{align*}

Подставим эти выражения в формулу для $R_{n+1}$:
%
\begin{align*}
    R_{n+1}
    &\approx
    R_n
    +
    \frac{
        -f'(x^*)R_n
        +
        \dfrac12 f''(x^*)R_n^2
    }{
        f'(x^*)
        -
        f''(x^*)R_n
    }.
\end{align*}

Приведём два слагаемых к общему знаменателю:
%
\begin{align*}
    R_{n+1}
    &\approx
    \frac{
        R_n\bigl(f'(x^*)-f''(x^*)R_n\bigr)
        -
        f'(x^*)R_n
        +
        \dfrac12 f''(x^*)R_n^2
    }{
        f'(x^*)-f''(x^*)R_n
    }
    \\
    &=
    \frac{
        -\dfrac12 f''(x^*)R_n^2
    }{
        f'(x^*)-f''(x^*)R_n
    }.
\end{align*}

Вынесем $f'(x^*)$ из знаменателя:
%
\begin{equation*}
    R_{n+1}
    \approx
    -
    \frac12
    \frac{f''(x^*)}{f'(x^*)}
    \frac{R_n^2}
    {1-\dfrac{f''(x^*)}{f'(x^*)}R_n}.
\end{equation*}

Вблизи корня $|R_n|$ мало, поэтому
%
\begin{equation*}
    \left|
        \frac{f''(x^*)}{f'(x^*)}R_n
    \right|
    \ll1.
\end{equation*}

Следовательно,
%
\begin{equation*}
    R_{n+1}
    \approx
    -
    \frac12
    \frac{f''(x^*)}{f'(x^*)}
    R_n^2.
\end{equation*}

Для абсолютной погрешности
%
\begin{equation*}
    \delta_n=|R_n|
\end{equation*}
%
получаем
%
\begin{equation*}
    \delta_{n+1}
    \approx
    \frac12
    \left|
        \frac{f''(x^*)}{f'(x^*)}
    \right|
    \delta_n^2.
\end{equation*}

Таким образом, погрешность следующего приближения пропорциональна
квадрату предыдущей погрешности:
%
\begin{equation*}
    \delta_{n+1}\sim\delta_n^2.
\end{equation*}

Поэтому метод Ньютона обладает \rtxt{3}{квадратичной сходимостью}.
При достаточно хорошем начальном приближении число верных знаков
примерно удваивается после каждой итерации.


\subsection{Недостатки метода}

Метод Ньютона очень быстро сходится вблизи простого корня, но имеет
два существенных недостатка.

Во-первых, он чувствителен к начальному приближению. Если $x_0$
выбрано вне области притяжения искомого корня, процесс может
расходиться либо сойтись к другому корню.

Во-вторых, на каждом шаге необходимо вычислять производную
%
\begin{equation*}
    f'(x_n).
\end{equation*}

Последний недостаток приводит к методу секущих.


% ============================================================
\section{Метод секущих}
% ============================================================

В методе Ньютона заменим производную приближённым значением,
полученным по двум последним точкам:
%
\begin{equation*}
    f'(x_n)
    \approx
    \frac{
        f(x_n)-f(x_{n-1})
    }{
        x_n-x_{n-1}
    }.
\end{equation*}

Подставляя эту оценку в формулу Ньютона,
%
\begin{align*}
    x_{n+1}
    &=
    x_n
    -
    \frac{f(x_n)}{f'(x_n)}
    \\
    &\approx
    x_n
    -
    \frac{
        f(x_n)
    }{
        \dfrac{f(x_n)-f(x_{n-1})}{x_n-x_{n-1}}
    }
    \\
    &=
    x_n
    -
    f(x_n)
    \frac{x_n-x_{n-1}}
    {f(x_n)-f(x_{n-1})}.
\end{align*}

Итак,
%
\begin{equation*}
    x_{n+1}
    =
    x_n
    -
    f(x_n)
    \frac{x_n-x_{n-1}}
    {f(x_n)-f(x_{n-1})}.
\end{equation*}

Этот процесс является двухшаговым: для нахождения $x_{n+1}$
необходимо знать два предыдущих приближения $x_{n-1}$ и $x_n$.

Геометрически вместо касательной используется секущая, проведённая
через точки
%
\begin{equation*}
    (x_{n-1},f(x_{n-1}))
    \qquad\text{и}\qquad
    (x_n,f(x_n)).
\end{equation*}
%
Её пересечение с осью $x$ задаёт следующее приближение.

\begin{center}
\begin{tikzpicture}[x=1.75cm,y=0.95cm]
    \pgfmathsetmacro{\xzero}{3.0}
    \pgfmathsetmacro{\xone}{2.35}
    \pgfmathsetmacro{\fzero}{ovff(\xzero)}
    \pgfmathsetmacro{\fone}{ovff(\xone)}
    \pgfmathsetmacro{\xtwo}{
        \xone-\fone*(\xone-\xzero)/(\fone-\fzero)
    }
    \pgfmathsetmacro{\ftwo}{ovff(\xtwo)}
    \pgfmathsetmacro{\xthree}{
        \xtwo-\ftwo*(\xtwo-\xone)/(\ftwo-\fone)
    }

    \draw[ovfaxis] (-0.10,0) -- (3.65,0)
        node[below,text=white] {$x$};
    \draw[ovfaxis] (0,-1.3) -- (0,3.65)
        node[left,text=white] {$y$};

    \draw[ovfcurve,domain=0.15:3.25,samples=120,smooth]
        plot (\x,{ovff(\x)});

    \draw[ovftangent]
        (\xzero,\fzero)
        --
        (\xtwo,0);

    \draw[ovftangent]
        (\xone,\fone)
        --
        (\xthree,0);

    \draw[ovfguide] (\xzero,0) -- (\xzero,\fzero);
    \draw[ovfguide] (\xone,0) -- (\xone,\fone);
    \draw[ovfguide] (\xtwo,0) -- (\xtwo,\ftwo);

    \fill[rc3] (\xzero,\fzero) circle (1.6pt);
    \fill[rc3] (\xone,\fone) circle (1.6pt);
    \fill[rc3] (\xtwo,\ftwo) circle (1.6pt);
    \fill[rc3] (1.2,0) circle (1.6pt);

    \node[text=white,below] at (\xzero,0) {$x_0$};
    \node[text=white,below] at (\xone,0) {$x_1$};
    \node[text=white,below] at (\xtwo,0) {$x_2$};
    \node[text=white,below] at (\xthree,0) {$x_3$};
    \node[text=white!85,below] at (1.2,0) {$x^*$};
    \node[text=rc4,anchor=south west] at (2.45,1.60) {$y=f(x)$};
\end{tikzpicture}

\captionof{figure}{Геометрическая схема метода секущих.}
\end{center}

Метод секущих не требует явного вычисления $f'(x)$ и обычно сходится
быстрее метода простых итераций, хотя порядок его сходимости ниже,
чем квадратичный порядок метода Ньютона.

\end{document}
